\preludeandfugue{Jennifer}
\contributor{Sadie Kaye}

% Timid Insecure Can't make up her mind, can't express herself

\prelude

Jennifer hated being called Jenny. That was what they had called 
her when she'd been a little girl. It was always `Jenny, you've 
got to do it this way,' and `Jenny, you know it's not safe to do 
that.' She'd seen the other girls and boys doing all sorts of 
things. They always seemed to have more fun than her, going out, 
playing together in the park or on the waste-ground that lay behind 
it. But she was never allowed to go out and play with them. `It 
wasn't safe,' they said.

She'd been a diligent but average schoolchild, doing the work that 
had been asked of her and rarely missing any assignment. Nevertheless, 
nothing she had learnt about really sparked anything inside her. 
She just got on with her work because they told her she must. And 
they had explained that if she did badly at her schoolwork and 
failed her exams it would make her life much more difficult later 
on. She believed them. Maths and English, they said, were 
particularly important. She worked hard at both. Spelling some 
words in English was difficult. So many words seemed to have many 
possible spellings and endings like --ant and --ent always sounded 
so similar that she never knew which was right and which was wrong. 
But when she did get it wrong she got a nasty red scrawl over her 
essay and she didn't like that at all. Maths too. It wasn't 
always easy to remember all those things they told her to do when 
they gave her maths work. And to be perfectly honest, she didn't 
always see why she had to do the calculations in the way they said. 
Some of the better maths-whizzes in the class said she was being 
silly and it was obvious what to do, but she didn't get it. Still, 
she hated the red crosses on her work and the `Jenny, you can do 
better than this,' comments. So, with a mixture of rote learning 
and hard work, she did well enough to pass the exams and get the 
sort of results that her parents were at least not very disappointed with.

It was important that her parents weren't cross with her for not 
being particularly good at maths or spelling, or any other subject. 
Fortunately, they never said much that was particularly critical. 
But they did go on rather about how clever her cousins were. They 
all had rather good results at A-level and were going to 
university.

Jennifer muddled through her A-levels. The subject she liked more 
than others was history. It always seemed to be a lot less 
effort to read a history textbook, absorb the facts she'd just read 
and write her history essay. Not everyone in her class said that 
remembering the details\=especially dates\=was straightforward, but 
it was not so difficult, she found. Her father, however, only had 
interests in science and technology, and dismissed any talk of 
history. Her mother, who had married him early without studying 
anything in particular, supported him. There was no point in 
studying the past, they said. What really matters is the here and 
now. And the future. How is the future and all the technology 
going to evolve? And in what way will you be part of it? There was 
so much happening in the world with the introduction of electronics 
and computers and silicon chips everywhere that it was important 
for a girl of Jenny's age to learn about these new innovations and 
not waste time with history and the past.

Jenny! No, she was Jennifer.

But in the end, she believed them. The future was where it was at. 
If only she'd known, she thought many years later, how difficult 
the future would be, and how important it was to strike out on 
one's own, but at the time her parents' advice seemed perfectly 
sensible.

She left school at eighteen and found ordinary work in a 
supermarket and later on waitressing at a local pizzeria. At some 
point only a couple of years later, she met Joe in a pub where she 
was with some people she knew. They chatted a bit. Or at least, he 
chatted to her and she listened hard and nodded. He had just 
returned from doing a degree at university and was starting up with 
a career in an insurance company. He seemed to like her and his 
career was progressing very well. In due course he proposed and 
they moved in together into their very own flat. To Jennifer, this 
was growing up at last. Adulting, as they call it. Her experience 
had been thankfully rather straightforward, and she looked forward 
to her future life as a housewife and hopefully as a mother.

\fugue

Jennifer was spending a rather normal day at home, doing the 
housework and planning the meals she'd be making in the next few 
days. It was at times like these, in the morning, when she was on 
her own in the house with nothing much to do\=nothing except to 
make sure everything was organised, clean and tidy\=that she 
thought about where she had come from and how she had got here. 
Joe's work had gone well and after only three years he had been 
earning enough for them to be able to buy their first house 
together. But the ghost of little Jenny followed her everywhere. 
Jenny who was always worried she'd do the wrong thing. Jenny who 
would be criticised for not managing her life properly. No, 
she told the ghost, emphatically. She'd really made it. There was nothing 
wrong. She had a secure life that was comfortable enough. But. But 
what? But nothing.

Joe, she knew, was sitting behind his desk at that very moment, 
with a computer in front of him and a telephone to the side. He'd 
be sending off emails and making phone calls and doing all his 
insurance work. She wasn't really sure what it exactly was that he 
did. He occasionally mentioned statistics and probability and used 
the complicated word `actuary' rather often, whatever that meant, 
but exactly what his work was and how it earned money was beyond 
her. But it did earn money and she was grateful for that.

Generally speaking, she kept the house clean and tidy. It was her 
work and she understood that she had to do her share of work too. 
And it was important that her work was good enough. Good enough 
for Joe when he came home and good enough for her parents and for 
Joe's parents when they came to visit. The ghost of Jenny said 
that she must avoid those carping teachers' reports\=even those 
from teachers from her own family after she'd grown up and become 
someone's wife. She must avoid all those ugly red Biro crosses on 
her work, avoid all the comments saying, `Jenny, you can do better 
than this.'

She did well. She did her work diligently with care. No one 
complained. No one gave her bad marks for housework. But there was 
one subject for which she only seemed to get bad marks, and she was 
afraid she would never pass the exams in. Oh, little Jenny was 
such a disappointment. It was true that this subject was hardly 
ever mentioned in conversation and certainly not raised directly to 
her\=not by Joe's parents nor by hers. Yet she could see in their 
faces all their disappointment and could hear their criticism in 
the way they skated round the subject, always avoiding it, always 
talking about other people or something else entirely.

Joe and she had been married five years now, but they had not yet 
had their first child.

As far as she could tell, this was perfectly fine for Joe. He had not been
complaining, but she felt the future grandparents' disappointment 
intensely. She didn't know why there was no child yet and she 
didn't know what do do. Her periods always arrived normally and she 
had stopped taking contraceptives when she and Joe had moved into 
the house together just over two years ago. But nothing had 
happened yet. It was not a subject she felt comfortable talking to 
Joe about. He seemed happy enough with the status quo. And wasn't 
her task in life to support him in what he wanted?

This was a frequent refrain going through her thoughts as she 
started with her daily housework. Perhaps that was because she 
always started with the bedroom where these thoughts were focused 
like a knot in her crotch and she only tidied and cleaned the other 
rooms later. And always, at this time in the morning, she wondered 
if she \textit{really} wanted children. She thought she did. When 
they had married she had been very sure of it. And, of course, as well as this,
she wanted what Joe wanted. At the same time, she wanted what 
her parents wanted and\=to only a slightly lesser extent\=what 
Joe's parents wanted. Yet no one ever talked about it. Still, she knew 
both parents hopes and wishes for her. She didn't have any 
brothers or sisters and Joe's only sibling was a sister in a 
relationship with another woman who didn't keep in touch with Joe 
or her parents-in-law except at Christmas and were showing no sign 
of any interest in children. Somehow, it seemed to be up to her 
and Joe.

It was still quite early in the morning. Her day always started 
with her making breakfast for Joe in her nightie with a bathrobe 
over. His breakfast was important to him, she knew. That was why it 
was her first task every day. When he had left for work she could 
look after herself a bit: this was the time for her to clean 
herself up and get dressed.

Today she was wearing one of her favourite dresses. It was nothing 
special, but it was perfect for the daytime and felt 
comfortable. She'd had it for only a few months but proportionately 
had worn it rather a lot since then. It was made of a thin and 
slightly stretchy dark green knitted fabric with a simple pattern 
showing white fern leaves. It fitted her well and, when she had 
bought it, Joe had said it showed off her figure nicely. She 
smiled to herself as she remembered him saying this\=the memory 
gave her a lot of pleasure. He didn't compliment her very often and 
as a result it was the more valuable when he did. She did have a 
good figure\=this was one of the things she could be proud of and 
another thing to keep an eye on and maintain. One more of her duties as 
a housewife. Her breasts weren't particularly large, but she made 
up for this with a flat tummy, fairly narrow waist and hips that 
were not so very broad. The dress was slim enough to show off her 
waist and stretchy enough to show off her modest breasts and hips. 
It still had on it a tag saying, `Do not remove this tag.' She 
didn't know what the tag was for, so she hadn't dared remove it.

She started in the bedroom, as was her habit. 
She looked round, scanned the room, putting away one or two 
items in drawers. This particular morning she noticed the bed's 
LEDs flashing.

As she watched the flashing lights, she recalled the excitement of 
moving into her first real house. Joe had insisted they do 
everything properly. They had enough money, he said, so they might 
as well. `Properly,' according to Joe, was to get the most modern 
of mod cons throughout. Internet of Things, he called it. 
Everything was connected to the internet. Even the bed. And AI 
was everywhere. Jennifer got out her phone and logged on to the 
bed's web page to see what it was that the bed wanted.

It turned out it was time to change the sheets. Somehow it felt to 
her like it was only the day before yesterday when she'd done that 
last, but Jennifer realised she had to obey. She started stripping 
the bed, putting aside the throw to go over the bed that did not 
need washing and also the duvet and pillows, but making a pile of 
the sheet, duvet cover and pillow cases. The bed's firmware kept a 
record of how often she did this. It could also inform her of the 
days when the mattress needed turning over or replacing completely. 
Fortunately, the current mattress, which had been new the day they 
moved in, would probably last another year.

She picked up the pile of dirty linen and took it to make a pile 
outside the washing machine and then returned to the bedroom. She'd 
deal with the washing later. Returning to the bedroom with clean 
linen from the cupboard on the landing, she made the bed carefully, 
making sure there were no creases in the sheets or lumps in the 
duvet. If there were any, it would be sure to let her know.

With the bed made, she went round the bedroom with duster and 
furniture polish. All the surfaces. Chest of drawers, wardrobe, 
dressing table. It was her usual routine, but she tried to avoid 
looking into the dressing table mirror. She really did not want to 
see what she looked like today.

Too late. It had seen her.

`You look tired,' said the dressing table. She kept silent but she 
felt its truth.

With an effort, she finished her dusting in the bedroom. Just as 
she'd finished and was leaving the bedroom to return the duster to 
the kitchen where it was kept, she passed the dressing table again.

`Jennifer?'

She turned and looked at it.

`Take ten minutes out for a cup of tea. Then come back and do the 
vacuuming.'

She nodded. She really needed a mid morning break.

With a comforting warm mug of tea in her hand she entered the living room 
and sat down. The sofa was on a side wall with the large window 
and glass door to the small verandah on her left and the doorway on 
her right. In front of her was a huge television. As she sat down 
with her tea, sensors in the sofa triggered the television on, and 
it started blurting out one of the many news channels.

`Volume down,' said Jennifer. She wasn't ready to hear whatever 
news there was today just yet and there was no way of turning the 
television off completely.

It was a lovely sunny day and the sky was a beautiful blue with a few 
cotton-wool clouds floating across. Jennifer realised she would 
much rather be sitting against the right hand wall of the sitting 
room, looking out of the big window at the sunshine outside. It 
would be easy enough to move the sofa to face a different 
direction. But as soon as she had had this thought she knew how 
impossible it was. The sofa had to face the television\=that was 
obvious, that was what Joe wanted. If she dared to move the sofa he 
might be tempted to move the television to some position in front 
of the window and that would block her view completely. She 
dismissed her alternative reality of enjoying sunshine from the 
window for what it was\=pure fantasy.

Suddenly the television channel switched over and the volume 
blurted back on loud.

`Switching to fantasy,' said the television. It was a film: one of 
the \textit{Lord of the Rings} films. `You'll like this one. It's 
the one where the dragon wakes up and starts roaring fire.'

`Turn it off.'

`You know I can't do that,' it said. 

`Volume to minimum then.'

With the television's sound still uncomfortably loud, Jennifer 
tried to ignore the television and focused on her mug of tea. If it 
wasn't for the hobbit trying to steal something from the dragon and 
the interminable background music she thought she might 
possibly be able to relax like this.

Ten minutes later\=no more, for the kettle in the kitchen where she 
had made her tea had been timing her break\=she got up from the 
sofa, got out the vacuum cleaner and took it into the bedroom. She 
vacuumed all round the bed and was just considering moving on to 
give the carpet on the landing her attention when her phone went 
off. It was some message. She turned off the vacuum cleaner and 
briefly enjoyed the sudden peaceful quiet this action brought. 
Then, remembering her phone had gone off, she wandered off to 
check.

It was a text from the hoover. The message said she'd missed a bit 
under the bed. There was an attachment with a photograph showing 
where she'd slipped up. Like a parking fine for housewives, she 
thought with a sarcastic smile on her face. Joe had recently 
received a parking fine. The letter had arrived in the post with a 
photograph of his car in the offending place where he had parked 
it. He had been furious.

Sighing a sigh, she returned, turned the vacuum cleaner back on, 
went to the dirty bit of carpet under the bed and made amends for 
her misdemeanour. With a sinking feeling,
she suddenly realised that she had no idea if the hoover's
email messages were all blind-copied to Joe at his work. And
all the other gadgets. Was that why he liked Internet of Things?

She tried to bury the thought and continued on through the rest of 
the house doing her best not to miss all those little out of the 
way spots where dust and dirt accumulate.

About an hour later, she had finished. She'd been everywhere with 
the vacuum cleaner, she thought. She turned it off and counted to 
a hundred, listening out for any further messages on her phone.

Nothing. She sighed a sigh of relief and enjoyed the peace for two 
whole minutes.

Then she carried the vacuum cleaner back to the cupboard where it 
was kept. One of the house pixies showed her where to stow it away. She 
said good morning to him and enquired after his health.

`Good, thank you,' he replied. `You did excellent work there. 
Don't forget that tomorrow you have to hoover the sofa, cushions 
and chairs.'

She nodded and turned to the kitchen.

She unstacked the dishwasher from the night before then cleared the 
breakfast table from Joe's breakfast. Another house pixie told her how to 
rinse the items and stack the dishwasher with the dirty crocks 
and cutlery, reminding her what items went where and how to avoid 
items touching each other so they do not scratch or break when the 
dishwasher is turned on. She knew all this of course\=she'd heard 
it hundreds of times\=but listened carefully and did what she had 
been told. Better safe than sorry, she told herself.

The dishwasher wasn't ready to go on yet. `I'm not ready to go on 
yet,' it said, out loud. So she moved on to filling the washing 
machine with the dirty sheets from the bed. Absentmindedly, she 
added detergent and twisted the knob in her usual fashion to select 
the programme, but the machine admonished her: `For that load the 
low temperature cotton programme is better.'

Silly Jenny!

No, she was Jennifer, and she was going to get it right. She 
changed the programme and turned the machine on.

Breathing deep breaths to keep calm she looked at the clock. It 
was lunch time. She'd not had any breakfast and suddenly realised 
that she was hungry.

As she made herself a simple lunch of pasta with pesto and a side 
salad, she thought about what she might make for Joe this evening 
for his main meal. She fancied making a fish pie for a change. 
There were some nice potatoes that would mash well and only 
yesterday she'd looked in the freezer and saw some tasty looking 
frozen cod that had been there for a while and needed using up. It 
had positively been begging to be cooked and eaten. And what's 
more, she would enjoy it too, for a change. Joe ate too much meat, 
she thought, and it would do them both good to have some fish for a 
change. And for pudding, an apple crumble would be perfect. It 
would be a chance both to finish off some apples that had been 
around a bit and for her to prove to Joe that she can cook, really. 
Not just ready meals, but her very own crumble topping and stewed 
apple. And the more she thought about it the more she realised she 
fancied cooking her own food tonight and making something she would 
like to eat as well, for a change. And she so much wanted Joe to know that she 
was capable of it.

`But \textit{you} know you \textit{can't} cook, not \textit{really},' replied the 
cooker with emphasis, reading her thoughts.

`Shut up cooker, what do \textit{you} know?'~said little Jenny, 
aged~14, entering from the hallway. She must have been playing 
Joe's computer game in his study. `You never ate any of Jennifer's 
meals.'

`I didn't have to. I heated them up and I inspected them. From 
what I saw, I really wouldn't have wanted to taste them even if I 
could.'

Jennifer buried her head in her hands as she sat at the table, her 
lunch mostly eaten. She was grateful to Jenny for supporting her 
but she just wanted some peace right now. Jenny often turned up 
this time of day. Sometimes the two of them had a long chat about 
the old days, but today she just wanted some quiet. In any case 
she had to prove to everyone that she really could cook.

She got out the kitchen scales, flour, sugar and butter. She might 
as well make a start and make the crumble topping. And the 
fish\=she should defrost that too. She went to the freezer to get 
the fish out.

`What are you doing? You know Joe hates your fish pies.' said the 
freezer. `It would be much better to make a shepherd's pie with 
the potatoes. Plenty of nice lamb mince here. Let me show you.'

The LEDs in the lamb mince started flashing red and green. Somehow 
the frozen cod she'd seen yesterday was nowhere to be found.

`Oh OK,' said Jennifer. `Since you insist.'

`And, while you're here,' said the freezer, `I really hope you are 
not planning to make an apple crumble for Joe. You know it will 
turn out badly and he will be unimpressed. There's a delicious 
ready-made apple pie here to defrost and cook; it's a new 
recipe and was a special at the supermarket. From the 
\textit{Chef's specials} range. Social media indicates a 98\% 
satisfaction rate for this pie. And what's more, house records 
show at least one tin of custard available. That way you won't 
have to impose on him that homemade slop you call custard either.'

`I suppose you are right,' said Jennifer and pulled out the mince 
and pie from the freezer.

Returning to the kitchen, she put away all the things for the 
crumble she'd planned. Little Jenny came up to her and gave her a 
hug.

`It's not as bad as all that,' she said. `You'll save a bit of 
time not making the crumble. Why don't you go for a walk in the 
park?'

Jennifer felt inadequate, but saw the sense in this. It was true, 
she realised. She had nearly two hours spare and a walk in the park 
was a lovely idea.

She checked out of the window. Bright sun, strong wind, not so 
warm.

`There could be a shower or two later on,' said little Jenny who was 
standing next to her.

Jennifer remembered how, one day coming home from school, she'd 
been caught in torrential rain. It had come from nowhere with no 
notice whatsoever. The cloud above burst open. In no time at all 
she had been soaked through. Soaked through her jacket to her 
shirt and down the inside of her skirt into her knickers. There 
was rain running down her legs into her shoes, soaking her socks. 
She walked home uncomfortably\=squelch squelch squelch\=as the sky 
brightened up. She was still soaked through when she arrived home
under a crystal blue sky.

When she got in the house, she burst into tears and her mother told 
her to `Snap out of it.' She went up to get into dry clothes but 
was never warm that evening. The next day she had a bad cold which 
she blamed on the rain, but her mother said there were no germs in 
rain. It was just an old wive's tale.

And now she was an old wife. Maybe not so old, but. What to wear 
outside if rain was possible? With her head low, she walked 
into her bedroom to get changed. There was an elderly woman sitting on 
a chair in the corner of the room.

`Who are you?'

`I'm Jennifer,' said the woman.

She turned her back to ignore the intruder, got out of one of the 
drawers a pair of jeans and put them on. The chest of drawers 
said, `Not a good idea. Cotton soaks up water and you will be cold 
and soaked through again.' She took them off and put on a pair of 
leggings. She felt cold.

Old Jennifer tut-tutted and asked what top she would wear. She 
hadn't decided. A non-cotton top layered-up with a fleece would be 
perfect, but she didn't have anything like this. In the end she 
screamed silently and put on an old T-shirt\=the only one she had 
that wasn't cotton\=and a jumper. She'd wear a raincoat over it all 
and people wouldn't see. And hopefully the rain\=if it 
came\=wouldn't be so strong to go through the coat and soak her 
through.

She inspected herself in the mirror and felt Little Jenny's eyes 
boring into her. `What are you going to wear over those 
leggings?'~she asked.

Jennifer had thought, `Maybe nothing,' but saw her point. Those legs. 
Especially the thighs. A miniskirt was ideal but the only ones she 
had were denim\=the evil cotton. A long woollen skirt? Scratches. 
Gets heavy when wet.

`You're over-thinking'~said the wardrobe. `If it doesn't rain 
you'll be boiled alive in all that.'

`And if it does rain hard you'll be soaked through just like 
before,' said Jenny.

`What will Joe think?'~said the old lady, watching avidly from the 
corner of the room.

`About the clothes I am wearing? That I am dowdy. It's his word. 
Sometimes he says I dress dowdy and I don't like it when he says 
that. But at least he can't see me now.'

`It works the other way too. What do you think he is doing right 
now?'

`He's at work,'~said Jennifer, and was suddenly unsure.

`If he falls out of love with you he has every opportunity to\='

`But he loves me. He said so last week.'

`Or \textit{when}, then. Not \textit{if}.'

`And what about your face?'~said the dressing table.

Jennifer looked at herself in the mirror. It had a point. There 
was a bad spot\=from an in-growing hair perhaps\=on her right cheek. 
She'd never be able to conceal that. And her eyebrows! When had 
she last plucked? Only last week, she thought, but it was a 
disaster. In desperation she grabbed the tweezers and started 
plucking. But her hurry made things worse rather than better.

`You are just like me, aren't you?'~said little Jenny, 
materialising from behind the dressing table mirror. `Silly Jenny, 
when you are not too timid you're always in a hurry. And the 
result is just as bad.'

She was right.

`What will Joe think?'~said the old lady.

There was no way Jennifer could go out today. She collasped on the 
bed and cried.

`Pull yourself out of that,' said the old teddy from the top shelf 
of the wardrobe. It was Jennifer's old teddy bear, she'd kept it 
as her only link to her childhood as Little Jenny. `Crying won't 
help anyone. Smarten yourself up and get ready for when Joe gets 
home.'

`At least someone has some sense round here,' said the old woman.

Jennifer nodded. She got off the bed, straightened the sheets and 
throw.

The teddy insisted on being taken down so he could see what was 
going on. Jennifer pulled him out and put him on the bedside 
table. She went back to the wardrobe and picked out one of her 
more revealing strappy dresses from the wardrobe. It was going to 
show a bit of cleavage, and she hoped she could get away with it.

`Oh, you're being brave now, wearing that,' said the teddy. `But 
will you pull it off?'

`More importantly,' said Old Jennifer, `will Joe?' She laughed, 
but Jennifer did not enjoy the joke.

There was a padded push-up bra that might help. She dug in the 
drawer for it, She took off her top and put the new bra on, and 
slipped the dress on over her head.

`That's better,' said Little Jenny.

She nodded and checked herself in the mirror. It was almost 
OK\=and would have to do. She flung off her trainers, leggings and 
socks and, still sobbing slightly, put on a pair of tights, picked 
out a pair of shoes and put them on.

`Make-up time!'~said the dressing table.

Struggling with her emotions, Jennifer attended to her face. 
Concealer to hide the spot and eyebrow pencil to tidy up the mess 
from earlier. Foundation. Powder. Eyeshadow, eyeliner\=this was 
always tricky\=and mascara. It would have to do.

`You mustn't cry any more. The make-up will run if you do,' said 
Old Jennifer.

Jennifer checked the time and wondered where it had all gone. She 
tidied up the clothes in the bedroom and then went off to the 
kitchen.

`Bring me!'~said the teddy. She picked him up and put him on the 
kitchen table. He was at least a memory of who she really was, or 
at least used to be.

She put on a pinafore, got on with peeling the potatoes, putting 
them on for mash, preparing carrots and broccoli and cutting up 
onions and putting togther the shepherd's pie. It was a meal she 
made frequently. One of the very few she could make that Joe 
seemed to like. Or at least, didn't complain about. Concentrating 
on her work seemed to help her focus and get all those voices out 
of her head.

`All except me,' said the teddy.

In due course she had the meal under control and had regained some 
of her composure. She hung up her pinafore and made another mug of 
tea but decided to drink it at the kitchen table to avoid the 
clairvoyant television. She was just draining her mug at six 
o'clock by the kitchen clock when Joe turned up, home from work.

`Hello darling,' said Jennifer. `Have you had a good day? It's
been very quiet here since you left this morning.'



